\documentclass[11pt]{article}
\usepackage[margin=1in]{geometry}
\usepackage{amsmath,amssymb}

\begin{document}
\begin{center}
{\large\bfseries When to compute the utility: a bound on the pair $(\mu_g,\sigma_g)$}
\end{center}
\medskip

\noindent
At a query image the GP gives a Gaussian posterior on the log firing rate,
$g\sim\mathcal{N}(\mu_g,\sigma_g^2)$, and the spike count in the window is Poisson
with rate $e^{g}$ (same setup as \texttt{laplace\_explainer.tex}). We want a rule
that decides, \emph{before} the utility is computed, whether the posterior at an
image is sane enough for that utility to mean anything. The rule is a bound on the
\emph{pair} $(\mu_g,\sigma_g)$, not on either one alone.

\section{A small $\sigma_g$ is not a small spread}
$\sigma_g$ is uncertainty in the \emph{log} rate, so it acts on the rate
multiplicatively: one step of $\sigma_g$ multiplies or divides the rate by
$e^{\sigma_g}$. A pair $(\mu_g,\sigma_g)$ that looks modest on paper can hide a
very skewed, very wide distribution of actual spike counts.

Take $\mu_g=3,\ \sigma_g=1.5$. The most probable count is just $2$ spikes and the
median is $20$ --- both look like an ordinary cell --- yet the mean is $62$, the
standard deviation $180$, and the $+3\sigma$ rate edge
$e^{\,\mu_g+3\sigma_g}=e^{7.5}\approx 1800$ spikes in the ${\sim}320$\,ms window
($\approx 5600$\,Hz, impossible). The most likely count and the mean differ
thirty-fold.

The three location measures of the count distribution $p(r)$ track those of the
underlying log-normal rate,
\[
\text{mode}\approx e^{\,\mu_g-\sigma_g^2},\qquad
\text{median}\approx e^{\,\mu_g},\qquad
\text{mean}=e^{\,\mu_g+\sigma_g^2/2},
\]
and they fan apart as $\sigma_g$ grows: the median stays pinned at the typical
rate $e^{\mu_g}$, while the mode slides down and the mean and the upper tail run
away.

\begin{center}
\small
\begin{tabular}{cccccccc}
\hline
$\mu_g$ & $\sigma_g$ & $\sigma_g^2$ & mode & median & mean & std & $e^{\mu_g+3\sigma_g}$\\
\hline
$1$ & $1.0$ & $1.00$ & $1$  & $3$  & $4$   & $6$   & $55$\\
$2$ & $1.0$ & $1.00$ & $3$  & $7$  & $12$  & $16$  & $148$\\
$3$ & $1.0$ & $1.00$ & $7$  & $20$ & $33$  & $44$  & $403$\\
$4$ & $1.0$ & $1.00$ & $20$ & $55$ & $90$  & $118$ & $1097$\\
\hline
$1$ & $1.5$ & $2.25$ & $0$  & $3$  & $8$   & $25$  & $245$\\
$2$ & $1.5$ & $2.25$ & $1$  & $7$  & $23$  & $66$  & $665$\\
$3$ & $1.5$ & $2.25$ & $2$  & $20$ & $62$  & $180$ & $1808$\\
$4$ & $1.5$ & $2.25$ & $6$  & $55$ & $168$ & $490$ & $4915$\\
\hline
\end{tabular}
\end{center}
\noindent
\textit{mode and median computed from $p(r)$ by deterministic quadrature; mean and
std are closed form; $e^{\mu_g+3\sigma_g}$ is the gate quantity of \S2.} Read
across a row: at fixed $\mu_g$, raising $\sigma_g$ from $1$ to $1.5$ barely moves
the median but multiplies the standard deviation and the $+3\sigma$ count
several-fold. A pair can look reasonable by its typical count and still encode a
distribution whose mean and upper tail are physically impossible.

\section{The gate}
The quantity that has to stay physical is the high end of the rate the posterior
still finds plausible, the $+3\sigma$ edge $e^{\,\mu_g+3\sigma_g}$. A retinal
ganglion cell produces at most of order a hundred spikes in the counting window,
so we require
\[
\boxed{\,e^{\,\mu_g+3\sigma_g}\;<\;T\,}
\qquad\Longleftrightarrow\qquad
\mu_g+3\sigma_g<\ln T .
\]
This couples the two parameters: a high mean leaves less room for the spread, and
vice versa. In the $(\sigma_g^2,\mu_g)$ plane the boundary
$\mu_g=\ln T-3\sqrt{\sigma_g^2}$ is a curve; in the $(\sigma_g,\mu_g)$ plane it is
the straight line $\mu_g=\ln T-3\sigma_g$. The multiplier $3$ is a margin choice:
it asks the $99.9$th rate percentile to stay below $T$; using $2$ would test only
the ${\sim}98$th, and Poisson noise pushes the true count a little past the rate
edge in any case.

\section{What $T$ should be: tie it to $r_{\max}$}
$T$ is not free. The utility truncates the Poisson sum at $r_{\max}$ (it sums over
$r=0,\dots,r_{\max}$), so the computed entropies are correct only if the count
distribution has negligible mass above $r_{\max}$. The edge
$e^{\,\mu_g+3\sigma_g}$ estimates where that mass ends, so the self-consistent gate
is
\[
e^{\,\mu_g+3\sigma_g}\;<\;r_{\max}.
\]
The live experiment used $r_{\max}=100$. With that truncation the honest threshold
is $T\approx 100$: above it the sum cannot see the counts the posterior places
there, and the utility is computed but wrong. Setting $T=1000$ is consistent only
if $r_{\max}$ is also raised to ${\approx}1000$. The clean choice is to fix
$r_{\max}$ from physics --- a few hundred counts, a few hundred Hz --- and set
$T=r_{\max}$. The gate and the truncation move together.

\section{A second reason, independent of the counts}
Even where the counts are physical, a large $\sigma_g$ degrades the Laplace
approximation the utility is built on (\texttt{laplace\_explainer.tex}): a small
$\sigma_g$ pins the integrand peak and the approximation is near-exact, a large
$\sigma_g$ widens it and the error grows. So past the gate the utility is suspect
twice over --- the distribution is unphysical \emph{and} the number computed from
it is a poor approximation.

\medskip
\noindent\textit{A flat cap $\sigma_g\le 1.5$ is the single-number shortcut for
this gate; it is conservative at low $\mu_g$ and too loose near $\mu_g\approx 4$,
because the real condition is on the pair $\mu_g+3\sigma_g$.}

\end{document}
